% ---------------------------------------------------------------------------
%  Recursive minimum-node AVL ("Fibonacci") tree macro.
%
%  T_h is drawn from T_{h-1} and T_{h-2} by the SAME recursion the proof uses,
%  and the horizontal layout is computed from phi via Binet's formula - so the
%  golden ratio literally positions the picture that proves the golden ratio.
%
%  Leaf count of T_h is L(h) = F(h+1); subtree centres are offset by half the
%  sibling's leaf count, which packs the tree with no overlap at any height.
%
%  Usage inside a tikzpicture:
%      \setavlscale{<leaf sep>}{<level sep>}
%      \drawfib{7}{0}{0}
% ---------------------------------------------------------------------------

\newcounter{avlnode}
\newcommand{\avlleafsep}{0.55}
\newcommand{\avllevsep}{0.80}

\newcommand{\setavlscale}[2]{%
  \renewcommand{\avlleafsep}{#1}%
  \renewcommand{\avllevsep}{#2}%
}

% Leaves of T_h via Binet: L(h) = F(h+1). Returns 0 for h < 0 (empty subtree).
\newcommand{\avlL}[1]{round((((1+sqrt(5))/2)^((#1)+1))/sqrt(5))}

% Node styles (overridden by the theme; these are safe standalone defaults)
\tikzset{
  avldot/.style   = {circle, draw=black, fill=black, inner sep=0pt, minimum size=3pt},
  avlgolddot/.style={circle, draw=black, fill=black, inner sep=0pt, minimum size=3pt},
  avledge/.style  = {draw=black, line width=0.4pt},
}

% \drawfib{height}{x}{y}
% Wrapped in a group so each frame's \avlself survives its own recursive calls;
% the counter is global, so node names stay unique across the whole picture.
\newcommand{\drawfib}[3]{%
  \begingroup
  \stepcounter{avlnode}%
  \edef\avlself{f\theavlnode}%
  \node[avldot] (\avlself) at (#2,#3) {};%
  \ifnum#1>0\relax
    % ---- left child is T(h-1); offset by half the RIGHT sibling's leaves ----
    \pgfmathsetmacro{\avlRsib}{\avlL{#1-2}}%
    \pgfmathsetmacro{\avlcxA}{(#2)-(\avlRsib/2)*\avlleafsep}%
    \pgfmathsetmacro{\avlcyA}{(#3)-\avllevsep}%
    \edef\avlnexth{\the\numexpr#1-1\relax}%
    \edef\avlchildA{f\the\numexpr\value{avlnode}+1\relax}%
    % fully expand the coordinates INTO the call, otherwise the recursive
    % invocation would re-bind \avlcxA before it is ever used
    \edef\avlgo{\noexpand\drawfib{\avlnexth}{\avlcxA}{\avlcyA}}\avlgo
    \draw[avledge] (\avlself) -- (\avlchildA);%
  \fi
  \ifnum#1>1\relax
    % ---- right child is T(h-2); offset by half the LEFT sibling's leaves ----
    \pgfmathsetmacro{\avlLsib}{\avlL{#1-1}}%
    \pgfmathsetmacro{\avlcxB}{(#2)+(\avlLsib/2)*\avlleafsep}%
    \pgfmathsetmacro{\avlcyB}{(#3)-\avllevsep}%
    \edef\avlnexthb{\the\numexpr#1-2\relax}%
    \edef\avlchildB{f\the\numexpr\value{avlnode}+1\relax}%
    \edef\avlgob{\noexpand\drawfib{\avlnexthb}{\avlcxB}{\avlcyB}}\avlgob
    \draw[avledge] (\avlself) -- (\avlchildB);%
  \fi
  \endgroup
}

% Convenience: reset the id counter before each tree so names are predictable
\newcommand{\resetavl}{\setcounter{avlnode}{0}}
